\chapter{\ifenglish Introduction\else บทนำ\fi}

\section{\ifenglish Project rationale\else ที่มาของโครงงาน\fi}
\enskip ในปัจจุบัน เทคโนโลยีปัญญาประดิษฐ์โดยเฉพาะด้านการประมวลผลภาพทางการแพทย์ ได้เข้ามามีบทบาทสำคัญในการช่วยวิเคราะห์และวินิจฉัยโรคจากภาพถ่ายทางคลินิก หนึ่งในงานที่ได้รับความสนใจคือการตรวจจับและแบ่งส่วนรอยโรคในภาพช่องปาก ซึ่งสามารถช่วยแพทย์ในการคัดกรองและติดตามอาการของผู้ป่วยได้อย่างมีประสิทธิภาพ

อย่างไรก็ตาม นอกจากการตรวจจับรอยโรคแล้ว การระบุว่า “รอยโรคเกิดขึ้นบนอวัยวะใด” เช่น กระพุ้งแก้ม ลิ้น หรือเพดานปาก ถือเป็นข้อมูลที่มีความสำคัญต่อการประเมินความรุนแรงของโรคและการวางแผนการรักษา แต่ในทางปฏิบัติ ภาพที่มีรอยโรคมักส่งผลให้ลักษณะทางกายวิภาคของอวัยวะผิดเพี้ยนไป ทำให้แบบจำลองการแบ่งส่วนอวัยวะทำงานได้ยากและมีความคลาดเคลื่อนมากขึ้น โดยเฉพาะในกรณีของภาพที่ไม่เคยปรากฏในชุดข้อมูลฝึก

จากปัญหาดังกล่าว งานวิจัยนี้จึงมุ่งเน้นการศึกษาความสามารถของแบบจำลองในการแบ่งส่วนอวัยวะในภาพช่องปากที่มีรอยโรคร่วมอยู่ รวมถึงการวิเคราะห์ว่าการประยุกต์ใช้กระบวนการ preprocessing และ post-processing สามารถช่วยเพิ่มประสิทธิภาพของแบบจำลองได้หรือไม่ เพื่อให้เข้าใจข้อจำกัดและศักยภาพของระบบในสถานการณ์จริงมากยิ่งขึ้น

\section{\ifenglish Objectives\else วัตถุประสงค์ของโครงงาน\fi}
\begin{enumerate}
    \item เพื่อศึกษาความสามารถของแบบจำลองการแบ่งส่วนอวัยวะ เมื่อทำการฝึกแบบจำลองโดยใช้ภาพที่มีรอยโรคร่วมอยู่
    \item เพื่อเปรียบเทียบประสิทธิภาพของแบบจำลองภายใต้กระบวนการประมวลผลที่แตกต่างกัน ได้แก่
    \begin{itemize}
        \item แบบจำลองพื้นฐาน (baseline model)
        \item แบบจำลองร่วมกับ preprocessing
        \item แบบจำลองร่วมกับ preprocessing และ post-processing
    \end{itemize}
\end{enumerate}

\section{\ifenglish Project scope\else ขอบเขตของโครงงาน\fi}

\subsection{\ifenglish Hardware scope\else ขอบเขตด้านฮาร์ดแวร์\fi}
การดำเนินงานวิจัยนี้ใช้ทรัพยากรด้านฮาร์ดแวร์ที่มีอยู่ ได้แก่ เครื่องคอมพิวเตอร์ส่วนบุคคลที่มีหน่วยประมวลผลกลาง (CPU) และหน่วยประมวลผลกราฟิก (GPU) สำหรับการฝึกและทดสอบแบบจำลอง โดยข้อจำกัดด้านฮาร์ดแวร์อาจส่งผลต่อขนาดของแบบจำลอง ระยะเวลาในการฝึก และปริมาณข้อมูลที่สามารถนำมาใช้ได้
\subsection{\ifenglish Software scope\else ขอบเขตด้านซอฟต์แวร์\fi}
โครงงานนี้พัฒนาและทดสอบแบบจำลองโดยใช้ภาษา Python ร่วมกับไลบรารีด้านการประมวลผลภาพและการเรียนรู้ของเครื่อง เช่น PyTorch, OpenCV และ NumPy รวมถึงเครื่องมือสำหรับการจัดการและวิเคราะห์ข้อมูล โดยไม่ได้มีการพัฒนา framework ใหม่ขึ้นมา
\subsection{ขอบเขตด้านการศึกษา}
\begin{enumerate}
    \item ศึกษาและประเมินแบบจำลองการแบ่งส่วนภาพ (image segmentation) สำหรับการระบุอวัยวะในภาพช่องปาก โดยใช้ข้อมูลภาพที่มีรอยโรคร่วมอยู่
    \item ใช้แบบจำลองที่มีอยู่แล้ว (pre-existing model) โดยไม่เน้นการพัฒนาโครงสร้างโมเดลใหม่ แต่เน้นการประเมินประสิทธิภาพของโมเดลภายใต้เงื่อนไขที่แตกต่างกัน
    \item ศึกษาผลของกระบวนการ preprocessing เช่น การปรับปรุงคุณภาพภาพ หรือการเตรียมข้อมูล ก่อนนำเข้าสู่แบบจำลอง
    \item ศึกษาผลของกระบวนการ post-processing เช่น rule-based methods หรือการปรับปรุงผลลัพธ์ segmentation เพื่อเพิ่มความถูกต้องของผลลัพธ์
    \item เปรียบเทียบประสิทธิภาพของระบบใน 3 รูปแบบ ได้แก่
        \begin{itemize}
            \item แบบจำลองพื้นฐาน (baseline)
            \item แบบจำลองร่วมกับ preprocessing
            \item แบบจำลองร่วมกับ preprocessing และ post-processing
        \end{itemize}
    \item การประเมินผลจะพิจารณาจากตัวชี้วัดด้านความแม่นยำของการแบ่งส่วน และการวิเคราะห์ข้อผิดพลาดของโมเดล โดยเฉพาะในภาพที่หรือไม่เคยปรากฏในชุดข้อมูลฝึก
\end{enumerate}

\section{\ifenglish Expected outcomes\else ประโยชน์ที่ได้รับ\fi}
\begin{enumerate}
    \item ทำให้ทราบถึงความสามารถและข้อจำกัดของแบบจำลองในการแบ่งส่วนอวัยวะในภาพช่องปากที่มีรอยโรค ซึ่งใกล้เคียงกับสถานการณ์ใช้งานจริง
    \item ช่วยให้เข้าใจผลกระทบของรอยโรคต่อความแม่นยำของการแบ่งส่วนอวัยวะ และลักษณะของข้อผิดพลาดที่เกิดขึ้น
    \item แสดงให้เห็นว่าการประยุกต์ใช้ preprocessing และ post-processing สามารถช่วยปรับปรุงประสิทธิภาพของแบบจำลองได้ในระดับใด
    \item เป็นแนวทางในการพัฒนาระบบวิเคราะห์ภาพทางการแพทย์ในอนาคต โดยเฉพาะในงานที่ต้องจัดการกับข้อมูลที่มีความซับซ้อนสูง
    \item สามารถนำผลการศึกษาไปใช้เป็นข้อมูลอ้างอิงในการออกแบบหรือปรับปรุง pipeline สำหรับงาน segmentation ในบริบทอื่น ๆ ได้
\end{enumerate}

\section{\ifenglish Technology and tools\else เทคโนโลยีและเครื่องมือที่ใช้\fi}

\subsection{\ifenglish Hardware technology\else เทคโนโลยีด้านฮาร์ดแวร์\fi}

\subsection{\ifenglish Software technology\else เทคโนโลยีด้านซอฟต์แวร์\fi}

\section{\ifenglish Project plan\else แผนการดำเนินงาน\fi}

\begin{plan}{6}{2020}{2}{2021}
    \planitem{7}{2020}{8}{2020}{ศึกษาค้นคว้า}
    \planitem{8}{2020}{1}{2021}{ชิล}
    \planitem{2}{2021}{2}{2021}{เผา}
    \planitem{12}{2019}{1}{2022}{ทดสอบ}
\end{plan}

\section{\ifenglish Roles and responsibilities\else บทบาทและความรับผิดชอบ\fi}
อธิบายว่าในการทำงาน นศ. มีการกำหนดบทบาทและแบ่งหน้าที่งานอย่างไรในการทำงาน จำเป็นต้องใช้ความรู้ใดในการทำงานบ้าง

\section{\ifenglish%
Impacts of this project on society, health, safety, legal, and cultural issues
\else%
ผลกระทบด้านสังคม สุขภาพ ความปลอดภัย กฎหมาย และวัฒนธรรม
\fi}

แนวทางและโยชน์ในการประยุกต์ใช้งานโครงงานกับงานในด้านอื่นๆ รวมถึงผลกระทบในด้านสังคมและสิ่งแวดล้อมจากการใช้ความรู้ทางวิศวกรรมที่ได้
